\noindent
Large language models are increasingly used as ``judges'' that stand in for human annotators, but the field lacks
an evidence-based account of \emph{when} an LLM judge agrees with humans well enough to substitute for them, and
\emph{what} governs that agreement. We address this for \emph{open, locally-servable} models, the realistic and
reproducible choice for most researchers. We contribute three studies anchored to a single principle: an LLM judge
should be measured against the \emph{human-human agreement ceiling} for a task, never a manufactured ``gold.''
\textbf{Study 1} is a PRISMA systematic review of \Nstudies{} studies across \Ndomains{} domains (all
\Nstudies{}${+}\STN$ references independently web-verified) that maps where LLM-as-a-judge agreement reaches the
human ceiling. \textbf{Study 2} replicates the standard summary-level G-Eval protocol on SummEval with the most
powerful open models and finds, with bootstrap confidence intervals, that open judges meet or exceed the established
proprietary-judge correlations on the \emph{verifiable} aspects (consistency, coherence) while the most compressed
aspect (fluency) is hard to rank for every judge. \textbf{Study 3} is a new meta-evaluation of a six-model open
panel (9B-754B) on natural-language-inference datasets with rich human label distributions ($\sim$100
annotators/item; \NLInjudg{} judgments), benchmarked against the annotator-bootstrap human ceiling with
chance-corrected agreement, jury aggregation, a calibrated contamination screen, and an item-level mixed-effects
analysis. The organizing finding is that, on ChaosNLI NLI, open judges \emph{track the human ceiling across the disagreement spectrum}: the
judge-vs-ceiling gap is near zero at both low and high human-entropy ($\NLIgapLow$, $\NLIgapHigh$), so accuracy
follows the ceiling rather than the judge failing on contested items; residual variance is item-level
($\mathrm{SD}_{\text{logit}}=\NLIitemSDlogit$ vs model $\NLImodelSDlogit$), and model scale contributes only modestly. We turn
this into a decision framework and an eight-point protocol for designing LLM-as-a-judge experiments. A companion
paper applies the same evaluation, with a systematic review, to subjective social-computing annotation in
Human-Centered Computing.
